Esta tarea fue desarrollada en el lenguaje C++. Cada paradigma algorítmico se encuentra separado en directorios distintos, facilitando su mantenimiento y evaluación independiente.

El algoritmo de fuerza bruta se encuentra en la carpeta code/brute\_force. El archivo principal es brute\_force.cpp, el cual lee los datos de entrada y procesa cada caso utilizando la función obtener\_diferencias, definida en algorithm/sequence\_difference.cpp. Esta función implementa una estrategia recursiva para generar todas las subsecuencias comunes posibles, seleccionando la más larga.

El algoritmo basado en programación dinámica se ubica en code/dynamic\_programming, con su archivo principal dynamic\_programming.cpp. Este programa también llama a obtener\_diferencias desde el subdirectorio algorithm, donde se construye la tabla de LCS y se reconstruye la diferencia entre las cadenas.

Ambos programas utilizan entradas desde archivos de texto ubicados en sus respectivas carpetas data/, y escriben la salida procesada y las mediciones en archivos separados. Cada uno cuenta con su propio makefile para compilar y ejecutar los casos de prueba automáticamente.

El código fuente completo, junto con los scripts de generación de entrada y graficación de resultados, se encuentra disponible en el siguiente repositorio privado:

\begin{mdframed}
    \begin{center}
        {\Large \url{https://github.com/VicenteVR/INF221-2025-1-TAREA-2-3}}
    \end{center}
\end{mdframed}
